\documentclass[conference]{IEEEtran}
\IEEEoverridecommandlockouts
% The preceding line is only needed to identify funding in the first footnote. If that is unneeded, please comment it out.
\usepackage{cite}
\usepackage{amsmath,amssymb,amsfonts}
\usepackage{algorithmic}
\usepackage{graphicx}
\usepackage{textcomp}
\usepackage{xcolor}
\def\BibTeX{{\rm B\kern-.05em{\sc i\kern-.025em b}\kern-.08em
    T\kern-.1667em\lower.7ex\hbox{E}\kern-.125emX}}
\begin{document}

\title{
\thanks{Identify applicable funding agency here. If none, delete this.}
}

\author{\IEEEauthorblockN{1\textsuperscript{st} Given Name Surname}
\IEEEauthorblockA{\textit{dept. name of organization (of Aff.)} \\
\textit{name of organization (of Aff.)}\\
City, Country \\
email address or ORCID}
\and
\IEEEauthorblockN{2\textsuperscript{nd} Given Name Surname}
\IEEEauthorblockA{\textit{dept. name of organization (of Aff.)} \\
\textit{name of organization (of Aff.)}\\
City, Country \\
email address or ORCID}}

\maketitle

\begin{abstract}
This survey paper provides a comprehensive overview of the rapidly evolving field of Semantic Simultaneous Localization and Mapping (SLAM), a crucial capability for mobile robots and autonomous systems. The integration of semantic information into traditional SLAM systems has emerged as a key theme, enabling improved localization accuracy, robustness, and mapping capabilities. Recent advances in visual SLAM, LiDAR SLAM, and point cloud processing have demonstrated the significance of incorporating semantic information to enhance system performance.

The paper reviews the fundamental concepts and techniques in Semantic SLAM, including the development of novel representations, such as Neural Radiance Fields, and the introduction of semantic-aided LiDAR SLAM systems. The importance of inertial measurement unit (IMU) integration, dynamic SLAM, and collaborative distributed SLAM systems is also highlighted. Furthermore, the paper discusses the role of deep learning in SLAM, real-time and efficient SLAM systems, and evaluation metrics and benchmarking methodologies.

The analysis of relevant papers reveals significant contributions to the field, including improved localization accuracy, effective loop closure detection, and the construction of globally consistent semantic maps. The paper also identifies open challenges and future research directions, such as integrating semantic information into SLAM systems and developing robust evaluation methodologies. This survey provides a timely and comprehensive overview of the current state and prospective future of Semantic SLAM technology, highlighting its potential to facilitate complex tasks, including autonomous navigation, scene understanding, and human-robot interaction. By summarizing key themes and findings, this paper aims to inform and inspire further research in this rapidly evolving field.
\end{abstract}

\begin{IEEEkeywords}
large language models, multimodal learning, natural language processing, machine learning, artificial intelligence
\end{IEEEkeywords}


\section{Introduction to Simultaneous Localization and Mapping (SLAM)}
Simultaneous Localization and Mapping (SLAM) is a crucial capability for mobile robots, enabling them to construct a global model of their environment from local observations while simultaneously determining their location within that environment \cite{b1}. The significance of SLAM in robotics cannot be overstated, as it underpins various applications including autonomous robot navigation, augmented reality (AR), and virtual reality (VR) \cite{b2}. Recent advances in SLAM have focused on improving the expressive capacity of environmental models, incorporating symbolic representation, and enhancing inference algorithms for more robust and accurate mapping \cite{b3}. This has led to significant improvements in the field, with researchers exploring various sensors and sensing modalities, such as monocular, stereo, and RGB-D cameras, to develop more sophisticated SLAM systems \cite{b4}.

The integration of semantic information with visual SLAM has shown promising results in dynamic and complex environments, where traditional SLAM approaches often struggle \cite{b5}. By incorporating high-level object and plane landmarks, symbolic representation, and learned representations, researchers have been able to develop more accurate and efficient SLAM systems \cite{b6}. However, current SLAM systems still face challenges in these scenarios, and ongoing research aims to develop more robust and certifiable inference methods, as well as more comprehensive environmental representations that incorporate hierarchical organization, affordance, dynamics, and semantics \cite{b7}. The pursuit of more comprehensive environmental representations is a key theme in SLAM research, with researchers seeking to model properties such as object relationships, scene understanding, and activity recognition \cite{b8}.

The development of novel classes of certifiable and robust inference methods has improved the reliability of SLAM systems in real-world operation \cite{b9}. Graphical models, machine learning, and optimization techniques have been used to improve SLAM performance, with researchers exploring various evaluation metrics and datasets to assess system performance \cite{b10}. Despite these advances, limitations and challenges remain, including the need for more robust data association, object recognition, and scene understanding in dynamic environments \cite{b11}. Furthermore, the integration of semantic information with visual SLAM requires careful consideration of issues such as uncertainty modeling, occlusion handling, and scalability \cite{b12}.

In conclusion, SLAM is a vital capability for mobile robots, with significant implications for various applications in robotics and beyond. Ongoing research in SLAM is focused on developing more robust, accurate, and comprehensive environmental representations, with a growing emphasis on the integration of semantic information and the development of certifiable inference methods \cite{b13}. As the field continues to evolve, we can expect to see significant advances in SLAM capabilities, enabling mobile robots to operate effectively in increasingly complex and dynamic environments. The connections between SLAM and other areas of research, such as computer vision, machine learning, and human-robot interaction, will be crucial in driving these advancements forward \cite{b14}. Ultimately, the development of more sophisticated SLAM systems will have far-reaching implications for fields such as robotics, AR/VR, and autonomous systems, enabling new applications and use cases that were previously unimaginable \cite{b15}.

\section{Background and History of SLAM: Evolution and Key Milestones}
The field of Simultaneous Localization and Mapping (SLAM) has undergone significant evolution over the past few decades, with various milestones marking its progress. The integration of semantic information into SLAM systems has emerged as a key theme, with studies demonstrating improved localization accuracy and loop closure detection through the use of semantic-aided LiDAR SLAM systems, such as SA-LOAM \cite{b1}. Furthermore, the development of new representations, including Neural Radiance Fields (NeRFs) and 3D Gaussian Splatting (3DGS), is reshaping the field of SLAM, enabling more accurate and efficient mapping \cite{b2}. These advancements have been facilitated by the creation of new datasets, such as DISCOMAN, which provide a benchmark for training and evaluating semantic SLAM methods \cite{b3}.

A comparative design space exploration of dense and semi-dense SLAM pipelines has yielded valuable insights into the algorithmic and hardware design spaces, highlighting the importance of robust benchmarking methodologies \cite{b4}. The need for such methodologies is emphasized by the lack of standard evaluation frameworks, which hinders the comparison of different SLAM systems. In contrast, some studies have focused on exploring different sensor combinations, such as the use of LiDAR and other sensors, to improve SLAM performance \cite{b5}. The choice of mapping representation also varies across studies, with some employing geometric representations and others utilizing semantic and radiance field-based representations.

The technical details and methodologies employed in SLAM systems are diverse, ranging from the analysis of KinectFusion and LSD-SLAM pipelines to the application of NeRFs for more accurate and efficient mapping \cite{b2}. Semantic segmentation has also become a crucial aspect of SLAM, with datasets like DISCOMAN providing ground truth occupancy grids and panoptic segmentation labels to facilitate the development of semantic SLAM methods \cite{b3}. Despite these advancements, several limitations and challenges persist in the field of SLAM, including loop closure detection, long-term autonomy, and the need for robust benchmarking methodologies \cite{b1} \cite{b5} \cite{b4}.

The evolution of SLAM has significant implications for various applications, including robotics, autonomous vehicles, and augmented reality. The integration of semantic information and the development of new representations have the potential to improve the accuracy and efficiency of SLAM systems, enabling more robust and reliable performance in complex environments. Furthermore, the creation of benchmarking datasets and frameworks will facilitate the evaluation and comparison of different SLAM systems, driving progress towards more accurate, efficient, and robust mapping and localization systems. By understanding these key findings, common themes, and technical details, researchers can build upon existing work and address the limitations and challenges in the field of SLAM, ultimately advancing the development of semantic SLAM and its applications.

\section{Fundamental Concepts and Techniques in SLAM: Overview and Analysis}
The field of Semantic SLAM has witnessed significant advancements in recent years, driven by the integration of semantic information into traditional SLAM systems. A thorough analysis of relevant papers reveals several key findings and contributions that have shaped the landscape of this research area. For instance, the introduction of a novel semantic-aided LiDAR SLAM system with loop closure detection, as proposed in \cite{b1}, has demonstrated improved localization accuracy and global consistent semantic mapping. Similarly, the work presented in \cite{b2} proposes a new SLAM system that utilizes semantic segmentation of objects and structures in the scene, resulting in a semantic map composed of clusters of points corresponding to object instances and structures.

The incorporation of deep learning techniques has also been a significant theme in Semantic SLAM research. The development of a deep semantic descriptor, as described in \cite{b3}, integrates semantic information into the feature extractor, thereby enhancing the robustness of SLAM methods. Furthermore, the use of decision trees to identify challenging benchmark properties and important components within the SLAM pipeline, as advocated in \cite{b4}, has highlighted the need for efficient and robust SLAM systems that can handle large-scale scenes and varying lighting conditions.

A common thread throughout these papers is the emphasis on the importance of semantic information in improving the accuracy, robustness, and efficiency of SLAM systems. The use of deep learning techniques, such as multi-task learning and shared encoder architectures, has been proposed as a means to learn semantic information and improve feature extraction \cite{b5}. Additionally, the need for efficient and robust SLAM systems that can handle loop closures and varying environmental conditions has been consistently highlighted \cite{b6}.

While there are no direct contradictions between the papers, some differing approaches and viewpoints are notable. For example, some papers focus on LiDAR-based SLAM \cite{b1}, while others consider camera-based SLAM \cite{b2}, \cite{b3}. The level of semantic information used also varies across papers, ranging from object detection and semantic segmentation to high-level scene understanding \cite{b7}. These differences in approach underscore the diversity of research directions in Semantic SLAM and highlight the need for further exploration and development.

The technical details and methodologies employed in these papers are also noteworthy. The use of SLAM frameworks, such as LOAM \cite{b1} and ORB-SLAM2 \cite{b2}, provides a foundation for the development of semantic SLAM systems. Semantic segmentation, as used in \cite{b2}, has been shown to be effective in creating semantic maps, while the integration of semantic information into feature extractors, as proposed in \cite{b3}, has improved the robustness of SLAM methods. Loop closure detection, a critical component of SLAM systems, has been addressed through the use of semantic graph-based place recognition methods \cite{b1}.

Despite the significant advancements in Semantic SLAM, several limitations and challenges remain. The need for efficient SLAM systems that can handle large-scale scenes and varying lighting conditions is a significant challenge \cite{b6}. Additionally, the robustness of SLAM systems to noise and variations in environmental factors, such as lighting and texture, remains an open issue \cite{b8}. Scalability is also a concern, as large-scale scenes and complex environments pose challenges for SLAM systems, requiring efficient and effective methods for loop closure detection and semantic mapping \cite{b9}.

In conclusion, the analysis of relevant papers in Semantic SLAM highlights the importance of incorporating semantic information, using deep learning techniques, and developing efficient and robust SLAM systems. The implications of these findings are significant, as they underscore the potential for Semantic SLAM to enable more accurate, robust, and efficient navigation and mapping in a wide range of applications, from robotics and autonomous vehicles to augmented reality and smart environments. As research in this area continues to evolve, it is likely that we will see further innovations and developments that address the limitations and challenges currently facing the field, ultimately leading to more effective and widely applicable Semantic SLAM systems \cite{b10}.

\section{Visual SLAM: Methods, Challenges, and Applications}
Visual SLAM: Methods, Challenges, and Applications

The development of visual Simultaneous Localization and Mapping (SLAM) systems has been a thriving area of research in recent years, with significant advancements in improving localization and mapping in dynamic and complex environments \cite{b1}. A key finding in this domain is the importance of incorporating semantic information into visual SLAM systems, which has shown promising results in enhancing the accuracy and robustness of these systems \cite{b1}, \cite{b3}, \cite{b5}. For instance, OpenVSLAM provides a versatile visual SLAM framework with high usability and extensibility, making it suitable for various applications such as augmented reality (AR) devices and autonomous control of robots and drones \cite{b2}.

The use of semantic information is a common theme among recent papers, with authors exploring its potential to improve visual SLAM systems. TextSLAM, for example, introduces a novel text-based visual SLAM approach that integrates planar text features into the visual SLAM pipeline, resulting in more robust and accurate 3D text maps \cite{b3}. Similarly, SGS-SLAM presents a semantic visual SLAM system based on Gaussian Splatting, addressing oversmoothing limitations of neural implicit SLAM systems and achieving state-of-the-art performance in camera pose estimation and map reconstruction \cite{b5}. These developments highlight the significance of semantic information in enhancing the performance of visual SLAM systems.

The application of visual SLAM in various fields such as AR/VR, autonomous robot navigation, and scene understanding is a recurring pattern in recent research \cite{b1}, \cite{b2}, \cite{b3}, \cite{b4}. However, despite these advancements, visual SLAM systems still face challenges in dynamic and complex environments, highlighting the need for further research and development \cite{b1}. The scalability and efficiency of visual SLAM frameworks remain a challenge, particularly in real-time applications \cite{b2}, \cite{b4}. Moreover, neural implicit SLAM systems are prone to oversmoothing limitations, which can be addressed using semantic Gaussian splatting approaches \cite{b5}.

Technical details and methodologies also play a crucial role in the development of visual SLAM systems. OpenVSLAM, for instance, provides a modular and extensible framework for visual SLAM, incorporating features such as illumination-invariant photometric error and semantic-guided keyframe selection \cite{b2}. TextSLAM represents text features compactly using three parameters and integrates them into the visual SLAM pipeline using illumination-invariant photometric error \cite{b3}. SGS-SLAM uses multi-channel optimization to incorporate appearance, geometry, and semantic features, addressing oversmoothing limitations of neural implicit SLAM systems \cite{b5}.

In conclusion, recent research in visual SLAM has highlighted the significance of incorporating semantic information into these systems, with various approaches and methodologies being proposed to address challenges such as dynamic and complex environments, scalability, and efficiency. As visual SLAM continues to evolve, it is likely that we will see further developments in semantic visual SLAM systems, with potential applications in a wide range of fields. The connections between visual SLAM and other areas of research, such as computer vision and robotics, will also be important to explore, as they have the potential to lead to significant advancements in our understanding of these systems and their applications \cite{b1}, \cite{b2}, \cite{b3}, \cite{b4}, \cite{b5}.

\section{Lidar SLAM and Point Cloud Processing: Advances and Limitations}
Lidar SLAM and point cloud processing have undergone significant advancements in recent years, driven by the increasing demand for accurate and efficient simultaneous localization and mapping (SLAM) systems. A key finding in this area is the importance of incorporating semantic information into Lidar SLAM and point cloud processing, as evident in the work presented by \cite{b1}, which introduces a novel semantic-aided LiDAR SLAM system, SA-LOAM. This approach leverages semantics in odometry and loop closure detection, resulting in improved localization accuracy and the construction of a globally consistent semantic map. Similarly, \cite{b2} proposes a fast algorithm for extracting semantically labeled surface segments from point clouds, enabling real-time processing and direct navigational use or higher-level contextual scene reconstruction.

The use of deep learning techniques has also become a prevalent theme in Lidar SLAM and point cloud processing, with several papers employing convolutional neural networks (CNNs) and self-supervised learning to improve semantic segmentation, keypoint detection, and description \cite{b3}, \cite{b4}. For instance, \cite{b5} presents a deep learned 3D keypoint detection and description method for LiDARs, exploiting intensity and depth information to extract powerful 3D features and enable robust scan-to-scan alignment and global localization. Furthermore, \cite{b6} proposes a self-supervised learning approach for semantic segmentation of lidar frames, training a deep point cloud segmentation architecture without human annotation and improving over time through multiple navigation sessions.

Despite these advancements, several challenges and limitations remain in Lidar SLAM and point cloud processing. One of the primary concerns is the need for efficient processing methods to handle large-scale point clouds and enable real-time applications \cite{b7}. Additionally, papers such as \cite{b8} highlight the difficulties posed by varying-density and occlusions in single-scan approaches, while \cite{b9} addresses the challenges of improving LiDAR semantic segmentation in long-distance outdoor scenarios. To mitigate these issues, researchers have proposed various technical details and methodologies, including point cloud processing techniques such as downsampling, plane constraint, and surface normal computation \cite{b10}, as well as voxel-based methods for improving LiDAR semantic segmentation \cite{b11}.

The implications of these developments are significant, with potential applications in fields such as robotics, autonomous vehicles, and augmented reality. The incorporation of semantic information into Lidar SLAM and point cloud processing has the potential to improve the accuracy and robustness of SLAM systems, while the use of deep learning techniques can enable more efficient and effective processing of large-scale point clouds. However, further research is needed to address the remaining challenges and limitations in this area, including the development of more efficient processing methods and the improvement of LiDAR semantic segmentation in long-distance scenarios.

In conclusion, Lidar SLAM and point cloud processing have undergone significant advancements in recent years, driven by the increasing demand for accurate and efficient SLAM systems. The incorporation of semantic information and the use of deep learning techniques have emerged as key themes in this area, with potential applications in a range of fields. However, further research is needed to address the remaining challenges and limitations, including the development of more efficient processing methods and the improvement of LiDAR semantic segmentation in long-distance scenarios. As researchers continue to explore and develop new techniques in this area, we can expect to see significant improvements in the accuracy, robustness, and efficiency of Lidar SLAM and point cloud processing systems \cite{b12}.

\section{Inertial Measurement Unit (IMU) Integration for Enhanced Accuracy}
Inertial Measurement Unit (IMU) Integration for Enhanced Accuracy

The integration of Inertial Measurement Unit (IMU) data with visual SLAM methods has emerged as a crucial aspect of Semantic SLAM, enabling significant improvements in accuracy and robustness \cite{b1}, \cite{b2}, \cite{b3}. The accurate preintegration of IMU measurements is a key factor in achieving good performance in deep inertial odometry \cite{b2}, and various approaches have been proposed to address this challenge. For instance, a novel dual visual-inertial SLAM network (DVI-SLAM) has been developed, which dynamically learns and adjusts the confidence maps of both visual factors and can be extended to include IMU factors \cite{b3}. The importance of integrating IMU data with visual SLAM methods is a common theme across various studies, highlighting the potential for improved performance and robustness in Semantic SLAM applications.

The use of deep learning architectures for inertial odometry and visual-inertial SLAM has become a prevalent approach \cite{b1}, \cite{b2}, \cite{b3}. These architectures have been shown to effectively leverage IMU data to enhance the accuracy and robustness of visual SLAM methods. Furthermore, benchmark datasets such as the TUM VI benchmark provide a comprehensive framework for evaluating visual-inertial odometry approaches \cite{b4}, enabling researchers to compare and refine their algorithms. The evaluation of visual-inertial odometry approaches using benchmark datasets has become a common practice, facilitating the development of more accurate and robust Semantic SLAM systems.

While some studies focus on the development of novel algorithms and architectures for visual-inertial SLAM \cite{b1}, \cite{b2}, \cite{b3}, others emphasize the importance of accurate preintegration of IMU measurements \cite{b2}, \cite{b5}. The "Trust Your IMU" study presents a contrasting viewpoint, arguing that modern pre-integration methods are accurate enough to ignore the drift for short time intervals \cite{b5}, whereas other studies may assume that IMU drift is a significant issue. This highlights the ongoing debate and research in this area, as scientists strive to optimize the integration of IMU data with visual SLAM methods.

The technical approaches employed in these studies include feature point matching algorithms optimized using IMU data \cite{b1}, deep learning architectures for inertial odometry and visual-inertial SLAM \cite{b2}, \cite{b3}, and accurate preintegration of IMU measurements using model-driven and data-driven approaches \cite{b2}, \cite{b5}. Experimental evaluations using real-world datasets \cite{b1}, \cite{b2}, \cite{b3}, \cite{b4} and comparisons with state-of-the-art algorithms and architectures \cite{b1}, \cite{b2}, \cite{b3}, \cite{b5} are also common methodologies used in these studies.

Despite the progress made in IMU integration for Semantic SLAM, several limitations and challenges remain. The need for accurate preintegration of IMU measurements to achieve good performance in deep inertial odometry \cite{b2} and the complexity of visual-inertial SLAM algorithms \cite{b1}, \cite{b3} are significant challenges. Additionally, the requirement for high-quality benchmark datasets to evaluate visual-inertial odometry approaches \cite{b4} and the assumption that IMU drift can be ignored for short time intervals may not always hold true in practice \cite{b5}. These limitations and challenges highlight the need for continued research and development in this area, as scientists strive to create more accurate and robust Semantic SLAM systems. Overall, the integration of IMU data with visual SLAM methods has the potential to significantly enhance the accuracy and robustness of Semantic SLAM applications, and ongoing research is focused on addressing the challenges and limitations associated with this approach.

\section{Dynamic SLAM: Handling Moving Objects and Scenes}
Dynamic SLAM: Handling Moving Objects and Scenes

The ability to handle dynamic scenes is crucial for Semantic SLAM systems, as real-world environments often feature moving objects and changing conditions \cite{b1}. Recent research has made significant progress in addressing this challenge, with various approaches focusing on identifying and removing dynamic objects or creating semantic maps that account for these changes. Det-SLAM, for example, combines ORB-SLAM3 and Detectron2 to achieve semantic SLAM in dynamic scenes by leveraging depth information and semantic segmentation \cite{b2}. Similarly, DynaSLAM II tightly integrates multi-object tracking and SLAM using instance semantic segmentation and ORB features, allowing it to optimize the structure of both static and dynamic components jointly \cite{b3}.

A common theme among these papers is the use of semantic segmentation to identify and remove dynamic objects or create semantic maps \cite{b2}, \cite{b3}, \cite{b5}. This approach enables systems like S3LAM to propose a new SLAM system that uses semantic segmentation of objects and structures, creating a semantic map and modifying the classical Bundle Adjustment formulation to constrain each cluster using geometrical priors \cite{b5}. Additionally, many papers leverage deep learning-based methods, such as Detectron2 and YOLO, for object detection and semantic segmentation, highlighting the importance of these techniques in handling dynamic scenes \cite{b2}, \cite{b4}.

The integration of multi-object tracking with SLAM is another significant development in Dynamic SLAM research. Papers like DynaSLAM II and Visual-Inertial Multi-Instance Dynamic SLAM demonstrate the effectiveness of tightly-coupled approaches, where object tracking and SLAM are optimized jointly \cite{b3}, \cite{b6}. In contrast, loosely-coupled approaches, such as Det-SLAM and Real-time SLAM Pipeline, also show promise in handling dynamic objects, albeit with a more modular design \cite{b2}, \cite{b4}.

The choice of camera system is another aspect that varies among the papers. While some use RGB-D cameras, like Real-time SLAM Pipeline, others employ stereo cameras, such as DynaSLAM II \cite{b3}, \cite{b4}. Moreover, the focus on static scene reconstruction differs between papers, with S3LAM emphasizing the creation of a semantic map and modification of Bundle Adjustment formulations, whereas other papers prioritize handling dynamic objects \cite{b5}.

Despite these advances, Dynamic SLAM still faces several challenges. Scalability and computational efficiency are significant concerns, as many systems struggle to handle complex scenes in real-time \cite{b2}, \cite{b3}. Robustness to varying lighting conditions is another challenge that researchers must address to ensure reliable performance in diverse environments \cite{b4}, \cite{b6}. Furthermore, integrating SLAM systems with other sensors, such as IMU, can potentially enhance their robustness and accuracy, as demonstrated by Visual-Inertial Multi-Instance Dynamic SLAM \cite{b6}.

In conclusion, recent research has made substantial progress in handling dynamic scenes in Semantic SLAM. By leveraging semantic segmentation, deep learning-based methods, and multi-object tracking, these systems have improved their ability to identify and remove dynamic objects or create semantic maps that account for changing conditions. However, challenges remain, and future work should focus on developing more efficient, scalable, and robust solutions that can handle complex scenes in real-time. As the field continues to evolve, we can expect to see further innovations and improvements in Dynamic SLAM systems, ultimately leading to more reliable and accurate Semantic SLAM applications \cite{b1}.

\section{Semantic SLAM: Incorporating High-Level Understanding into Mapping}
Semantic SLAM: Incorporating High-Level Understanding into Mapping

The development of Semantic SLAM has been rapidly advancing in recent years, with a growing focus on incorporating high-level understanding into mapping \cite{b1}. This paradigm shift is driven by the need for more accurate and informative maps that can facilitate complex tasks such as autonomous navigation, scene understanding, and human-robot interaction. Key findings from recent papers have contributed significantly to this field, including the introduction of hierarchical categorical representations for 3D Gaussian Splatting \cite{b2}, deep semantic descriptors learned through multi-task learning \cite{b3}, and semantic-aided LiDAR SLAM systems with loop closure detection \cite{b4}. These advancements have enabled accurate global 3D semantic mapping, improved feature extraction and matching accuracy, and enhanced localization accuracy in large-scale environments.

A common theme among these papers is the integration of semantics into SLAM systems to improve mapping accuracy and enable high-level understanding \cite{b5}. This is achieved through the use of deep learning techniques, such as convolutional neural networks (CNNs) and multi-task learning, which learn semantic representations and improve feature extraction \cite{b6}. The focus on large-scale environments is also a prominent trend, with several papers addressing the challenge of scaling up SLAM systems to handle complex indoor and outdoor spaces \cite{b7}. Furthermore, the importance of datasets and benchmarking is highlighted by the introduction of new datasets, such as DISCOMAN, which provides a comprehensive platform for training and evaluating semantic SLAM methods \cite{b8}.

Despite these advancements, there are contrasting viewpoints and approaches in the field. For example, Hi-SLAM proposes a hierarchical categorical representation, whereas other papers use flat semantic representations \cite{b9}. Additionally, some papers focus on LiDAR-based SLAM, while others use vision-based approaches \cite{b10}. The use of external data sources is also a point of contention, with some papers leveraging external data to enhance semantic mapping, while others rely solely on onboard sensors \cite{b11}.

The technical methodologies employed in these papers are diverse and include 3D Gaussian Splatting, multi-task learning, semantic-assisted ICP, and CNNs for semantic segmentation and feature extraction \cite{b12}. However, the limitations and challenges of semantic SLAM systems are also acknowledged, including scalability issues, data quality and availability, computational efficiency, and integration with other components \cite{b13}. These challenges highlight the need for further research and development to create more robust, efficient, and scalable semantic SLAM systems.

The implications of these developments are significant, as they have the potential to enable more accurate and informative mapping, facilitate complex tasks such as autonomous navigation, and enhance human-robot interaction \cite{b14}. The connections between semantic SLAM and other fields, such as computer vision, robotics, and artificial intelligence, are also noteworthy, as they highlight the interdisciplinary nature of this research area \cite{b15}. Overall, the field of semantic SLAM is rapidly advancing, with a growing focus on incorporating high-level understanding into mapping, and has significant implications for a wide range of applications.

\section{Collaborative and Distributed SLAM Systems for Multi-Robot Environments}
Collaborative and distributed SLAM systems for multi-robot environments have garnered significant attention in recent years, driven by the potential to enhance localization accuracy, reduce uncertainty, and improve mapping capabilities \cite{b1}. A key contribution in this area is the introduction of Swarm-SLAM, a sparse decentralized collaborative SLAM framework that supports various sensing modalities and inter-robot loop closure prioritization \cite{b2}. This approach enables robots to operate autonomously, reducing communication overhead and facilitating scalable solutions for large teams of agents. In contrast, A Collaborative Visual SLAM Framework presents an edge server-based approach for service robots, enabling real-time information sharing and landmark organization \cite{b3}. This centralized paradigm allows for efficient collaboration between robots, but may be limited by the need for a reliable communication infrastructure.

The importance of decentralization in collaborative SLAM systems is a common theme across many papers, with Kimera-Multi proposing a fully distributed multi-robot system for dense metric-semantic SLAM \cite{b4}. This approach builds 3D mesh models with semantic labels and leverages inter-robot loop closures to improve localization accuracy and reduce uncertainty. Similarly, COVINS introduces a centralized collaboration paradigm for visual-inertial SLAM, enabling scalable SLAM in large environments and teams of agents \cite{b5}. However, this approach relies on a server back-end to establish a global estimate, which may be limiting in scenarios where communication infrastructure is unreliable.

Inter-robot collaboration is another key aspect of collaborative SLAM systems, with all papers highlighting the benefits of cooperation between robots \cite{b1}, \cite{b2}, \cite{b3}, \cite{b4}, \cite{b5}. By sharing information and coordinating their actions, robots can improve their understanding of the environment and reduce the uncertainty associated with localization and mapping. Furthermore, semantic mapping is an important development in collaborative SLAM, providing a more comprehensive understanding of the environment and enabling robots to make informed decisions about their actions \cite{b4}.

The technical details of collaborative SLAM systems are also noteworthy, with Kimera-Multi introducing a novel incremental maximum clique outlier rejection for distributed place recognition and robust pose graph optimization \cite{b4}. Additionally, A Collaborative Visual SLAM Framework proposes an elegant communication pipeline and landmark organization method to enable real-time information sharing between robots \cite{b3}. Mesh deformation techniques are also used in Kimera-Multi to correct local meshes based on improved trajectory estimates \cite{b4}.

Despite the advancements in collaborative SLAM systems, there are still several limitations and challenges that need to be addressed. Communication overhead is a significant issue in decentralized approaches, which can be mitigated using techniques like inter-robot loop closure prioritization \cite{b2}. Scalability limitations are also a concern, with some papers demonstrating the ability to scale collaborative SLAM systems to large teams of agents or complex environments, while others acknowledge the challenges associated with this \cite{b1}, \cite{b5}. Furthermore, sensor noise and limitations can impact collaborative SLAM performance, highlighting the importance of understanding sensor characteristics and limitations for long-term autonomy \cite{b6}.

In conclusion, collaborative and distributed SLAM systems for multi-robot environments have made significant progress in recent years, with key developments including decentralization, inter-robot collaboration, and semantic mapping. While there are still several limitations and challenges that need to be addressed, the potential benefits of these systems make them an exciting area of research, with implications for a wide range of applications, from service robotics to autonomous vehicles \cite{b1}, \cite{b2}, \cite{b3}, \cite{b4}, \cite{b5}. As the field continues to evolve, it is likely that we will see the development of more sophisticated collaborative SLAM systems, capable of operating in complex environments and providing accurate and reliable localization and mapping capabilities.

\section{Deep Learning in SLAM: Recent Advances, Challenges, and Future Directions}
Deep learning has revolutionized the field of Simultaneous Localization and Mapping (SLAM), enabling the development of more robust, efficient, and scalable systems. Recent advances in deep learning-based SLAM have led to significant improvements in various applications, including autonomous vehicles, intelligent robots, and LiDAR-based navigation systems. This section provides an overview of the recent advances, challenges, and future directions in deep learning-based SLAM.

The integration of deep learning techniques with traditional SLAM systems has been a key theme in recent research. For instance, the use of differentiable SLAM architectures as a training signal for deep learning-based models has improved performance in various LiDAR-based applications, such as classification, regression, and SLAM \cite{b1}. Additionally, the use of deep local feature descriptors obtained by neural networks has improved the efficiency and stability of traditional SLAM systems \cite{b2}. The incorporation of semantic information into SLAM systems has also been shown to improve localization accuracy, detect loop closures effectively, and construct a global consistent semantic map \cite{b3}.

The use of deep learning techniques has also enabled the development of more robust and adaptable keypoint extraction methods. For example, the Model-Agnostic Meta-Learning (MAML) strategy has been used to optimize the training of keypoint extraction networks, enhancing their adaptability to diverse environments \cite{b4}. Furthermore, recent advancements in computer vision have provided a data-driven approach to tackle the Visual-SLAM problem, enabling the development of more efficient and scalable systems \cite{b5}.

Despite these advances, several challenges and limitations remain. Cumulative positioning errors in SLAM systems can be mitigated using online learning modules \cite{b4}, but generalization in continuous motion scenes remains a challenge, affecting loop detection accuracy. Moreover, the scalability and efficiency of deep learning-based SLAM systems are crucial for real-time operation, particularly in complex environments.

The papers analyzed in this section highlight the importance of semantic information in improving localization accuracy and detecting loop closures. The integration of deep learning techniques with traditional SLAM systems has also been shown to improve performance, efficiency, and robustness. However, contrasting viewpoints and approaches have been presented, including traditional vs. learning-based methods \cite{b5} and differentiable SLAM vs. traditional SLAM \cite{b1}. Technical details and methodologies employed in these papers include deep learning architectures, such as convolutional neural networks (CNNs) and recurrent neural networks (RNNs), and various SLAM systems, including ORB-SLAM3 and LIFT-SLAM.

In conclusion, deep learning has significantly advanced the field of SLAM, enabling the development of more robust, efficient, and scalable systems. However, challenges and limitations remain, including cumulative positioning errors, generalization in continuous motion scenes, and scalability and efficiency. Future research directions should focus on addressing these challenges and exploring new applications and techniques, such as the use of multimodal sensors and the integration of deep learning with other AI techniques. The implications of these developments are significant, with potential applications in various fields, including robotics, autonomous vehicles, and augmented reality. As the field continues to evolve, it is essential to consider the connections between deep learning, SLAM, and other areas of research, such as computer vision and machine learning, to fully realize the potential of these technologies.

\section{Real-Time and Efficient SLAM Systems for Practical Applications}
Real-Time and Efficient SLAM Systems for Practical Applications

The development of real-time and efficient Simultaneous Localization and Mapping (SLAM) systems is crucial for their practical applications in various fields, including robotics, autonomous vehicles, and augmented reality \cite{b1}. Recent research has focused on achieving real-time performance while maintaining accuracy and robustness. For instance, the work presented in \cite{b2} emphasizes the importance of systematic quantitative evaluation of SLAM algorithms to navigate the landscape for real-time localization and mapping. This approach enables the automated exploration of algorithmic and implementation design space, leading to accelerated adaptive SLAM solutions.

The incorporation of semantic information into SLAM systems has also been a key theme in recent research \cite{b3}, \cite{b4}. The SeMLaPS system, proposed in \cite{b3}, combines 2D neural networks and 3D networks based on SLAM systems with 3D occupancy mapping to achieve real-time semantic mapping from RGB-D sequences. Similarly, the SA-LOAM system presented in \cite{b4} leverages semantics in odometry and loop closure detection, resulting in improved localization accuracy and effective loop closure detection. These works demonstrate the benefits of integrating semantic information into SLAM systems for enhanced accuracy and robustness.

Efficient feature selection is another critical aspect of real-time SLAM systems \cite{b5}, \cite{b6}. The greedy-based feature selection method proposed in \cite{b5} significantly improves the accuracy and efficiency of an L-SLAM system by actively selecting a subset of features. This approach uses a stochastic-greedy algorithm to approximate optimal results in real-time, reducing computational complexity. Additionally, the iMAP system presented in \cite{b6} employs a multilayer perceptron (MLP) as the only scene representation, building a dense, scene-specific implicit 3D model of occupancy and color.

The use of hybrid approaches has also been explored in recent research \cite{b3}, \cite{b4}. The combination of different techniques, such as 2D and 3D networks or semantic graph-based place recognition, has shown promise in achieving real-time performance while maintaining accuracy. However, the evaluation of SLAM algorithms remains a crucial aspect, with some papers focusing on systematic quantitative methods \cite{b2} and others proposing new algorithms without extensive evaluation.

The technical details and methodologies employed in these works vary, with machine learning techniques such as latent prior networks \cite{b3}, stochastic-greedy algorithms \cite{b5}, and multilayer perceptrons \cite{b6} being used. The SLAM system architectures also differ, including 3D occupancy mapping \cite{b3}, LiDAR-based systems \cite{b4}, \cite{b5}, and RGB-D camera-based systems \cite{b6}. Despite these variations, the common themes of real-time processing, semantic mapping, efficient feature selection, and hybrid approaches highlight the current trends in SLAM research.

However, several limitations and challenges remain, including computational complexity, data quality and availability, and robustness and generalizability \cite{b1}-\cite{b6}. The need to reduce computational complexity while maintaining accuracy is a recurring theme, with many papers acknowledging the importance of efficient feature selection and real-time processing. Additionally, the reliance on high-quality data for training and operation can be a limitation in certain scenarios, highlighting the need for more robust and generalizable SLAM systems.

In conclusion, recent research has made significant progress in developing real-time and efficient SLAM systems for practical applications. The incorporation of semantic information, efficient feature selection, and hybrid approaches have been key themes, with machine learning techniques and varying SLAM system architectures being employed. However, limitations and challenges remain, emphasizing the need for continued research to achieve more robust, generalizable, and computationally efficient SLAM systems \cite{b1}-\cite{b6}.

\section{Evaluation Metrics and Benchmarking for SLAM Algorithms}
Evaluation Metrics and Benchmarking for SLAM Algorithms

The development of robust evaluation metrics and benchmarking methodologies is crucial for advancing the field of Semantic SLAM \cite{b1}. Recent papers have made significant contributions to this area, providing valuable insights into the characterization of SLAM benchmarks \cite{b2}, comparative design space exploration \cite{b3}, and the introduction of large-scale benchmark datasets \cite{b4}. These efforts highlight the importance of understanding what factors lead to track failure or degradation in state-of-the-art SLAM algorithms. For instance, a decision tree-based approach has been proposed to identify challenging benchmark properties \cite{b2}, while an extended SLAMBench framework has been used to compare two state-of-the-art SLAM pipelines, namely KinectFusion and LSD-SLAM \cite{b3}.

A common theme emerging from these papers is the need for robust benchmarking methodologies to evaluate SLAM algorithms holistically and quantitatively \cite{b1}. This is particularly important in Semantic SLAM, where the use of semantic information can significantly improve performance, especially in loop closure detection and odometry \cite{b5}. The development of flexible and modular frameworks, such as XRDSLAM \cite{b6}, has also been highlighted as essential for facilitating the combination of different algorithm modules and standardized benchmarking. Furthermore, the importance of profiling computational performance and optimizing SLAM components for efficient execution cannot be overstated \cite{b7}.

The papers also present some contrasting viewpoints or approaches, such as qualitative vs. quantitative evaluation \cite{b8} and geometry-based vs. learning-based SLAM \cite{b9}. While qualitative visualizations can provide valuable insights into SLAM performance, quantitative evaluation metrics are essential for comparing different algorithms and systems. Additionally, the comparison of classical geometry-based and recent learning-based SLAM algorithms has highlighted the strengths and weaknesses of each approach \cite{b9}. Different SLAM pipelines, such as KinectFusion and LSD-SLAM, have also been compared, demonstrating the importance of understanding their algorithmic and hardware design spaces \cite{b3}.

From a technical perspective, the papers employ a range of methodologies, including decision trees \cite{b2}, SLAMBench framework \cite{b3}, physically-based rendering \cite{b10}, semantic graph-based place recognition \cite{b11}, and modular code design and multi-process running mechanism \cite{b6}. These technical details and methodologies have significant implications for the development of robust evaluation metrics and benchmarking methodologies in Semantic SLAM. For example, the use of decision trees can help identify challenging benchmark properties, while the SLAMBench framework can facilitate the comparison of different SLAM pipelines.

Despite these advances, several limitations and challenges remain, including the lack of benchmarking methodologies \cite{b1}, computational efficiency \cite{b7}, generalization to unseen data \cite{b12}, and integration of different algorithm modules \cite{b6}. Addressing these challenges will be crucial for advancing the field of Semantic SLAM and developing robust evaluation metrics and benchmarking methodologies. The development of large-scale benchmark datasets, such as the one introduced in \cite{b4}, can help alleviate some of these challenges by providing a common platform for evaluating and comparing different SLAM algorithms and systems.

In conclusion, the development of robust evaluation metrics and benchmarking methodologies is essential for advancing the field of Semantic SLAM. Recent papers have made significant contributions to this area, highlighting the importance of understanding what factors lead to track failure or degradation in state-of-the-art SLAM algorithms. The use of semantic information, flexible and modular frameworks, and profiling computational performance are all crucial aspects of developing robust evaluation metrics and benchmarking methodologies. Addressing the limitations and challenges remaining in this area will be essential for developing robust and efficient Semantic SLAM systems \cite{b13}.

\section{Open Challenges and Future Research Directions in the Field of SLAM}
The field of Semantic SLAM has witnessed significant advancements in recent years, with various approaches and techniques being proposed to improve localization accuracy, robustness, and mapping capabilities \cite{b1}, \cite{b2}. However, despite these developments, several open challenges and future research directions remain to be addressed. One of the key themes that emerges from the analysis of relevant papers is the importance of integrating semantic information into SLAM systems \cite{b3}, \cite{b4}. This can be achieved through various means, such as leveraging convolutional neural networks (CNNs) for semantic information extraction and association \cite{b5}, or utilizing techniques like semantic-assisted Iterative Closest Point (ICP) algorithm for improved odometry and loop closure detection \cite{b1}.

Another significant challenge in the field of Semantic SLAM is the need for standardized benchmarking and evaluation frameworks \cite{b2}, \cite{b6}. This is crucial for effectively evaluating the performance of different SLAM systems and algorithms, as well as facilitating comparisons and identifying areas for improvement. The development of unified benchmarking frameworks, such as the one proposed in \cite{b2}, can help address this challenge and provide a common platform for researchers to test and validate their approaches.

Long-term autonomy is another critical aspect that requires attention in the context of Semantic SLAM \cite{b7}, \cite{b8}. As SLAM systems are expected to operate over extended periods, they must be able to adapt to changing environments, manage resources efficiently, and maintain robustness in the face of various challenges. This necessitates the development of advanced techniques for semantic information integration, sensor fusion, and scene understanding, as well as more sophisticated methods for loop closure detection, relocalization, and mapping \cite{b9}, \cite{b10}.

The survey of relevant papers also highlights the significant impact of advances in computer vision and machine learning on the development of SLAM systems \cite{b11}, \cite{b12}. Techniques like Neural Radiance Fields (NeRFs) and 3D Gaussian Splatting have shown great promise in improving scene representation, rendering, and understanding \cite{b13}, \cite{b14}. However, further research is needed to fully exploit the potential of these techniques in the context of Semantic SLAM and to address the challenges associated with their integration into existing SLAM frameworks.

Furthermore, the analysis of papers reveals that scalability and generalization remain significant challenges in the field of Semantic SLAM \cite{b15}, \cite{b16}. As SLAM systems are expected to operate in large, complex environments, they must be able to scale up efficiently and generalize well to unseen data. This requires the development of more advanced techniques for semantic information integration, feature extraction, and scene understanding, as well as more sophisticated methods for loop closure detection, relocalization, and mapping \cite{b17}, \cite{b18}.

In addition, the quality and availability of semantic information can significantly impact the performance of Semantic SLAM systems \cite{b19}, \cite{b20}. High-quality semantic information may not always be available or reliable, which can lead to errors in localization, mapping, and scene understanding. Therefore, further research is needed to develop more robust methods for semantic information extraction, association, and integration, as well as techniques for handling uncertain or incomplete semantic information \cite{b21}, \cite{b22}.

In conclusion, the field of Semantic SLAM faces several open challenges and future research directions, including the integration of semantic information, benchmarking and evaluation, long-term autonomy, advances in computer vision and machine learning, scalability and generalization, and semantic information quality and availability. Addressing these challenges will require significant advancements in various areas, including sensor fusion, scene understanding, loop closure detection, relocalization, and mapping. By exploring these research directions and developing more sophisticated techniques for Semantic SLAM, researchers can create more robust, efficient, and autonomous systems that can operate effectively in a wide range of environments and applications \cite{b23}, \cite{b24}.

\section{Conclusion: The Current State and Prospective Future of SLAM Technology}
The current state of SLAM technology, as reflected in the analyzed papers, underscores significant advancements in the field, particularly with the integration of semantic information into SLAM systems \cite{b1}. The development of benchmarking frameworks such as SLAMBench \cite{b2} and GSLAM \cite{b3} has been instrumental in enabling holistic and quantitative evaluation of SLAM systems, highlighting the importance of robust evaluation methodologies. Moreover, the incorporation of semantic information, as seen in SA-LOAM \cite{b4}, has improved localization accuracy and loop closure detection, paving the way for enhanced performance and high-level intelligence in SLAM systems.

A common theme emerging from the papers is the need for modular and flexible frameworks that facilitate the development of complete SLAM systems and standardized benchmarking \cite{b5}. XRDSLAM \cite{b6}, for instance, employs a modular code design and multi-process running mechanism to achieve this goal. The consideration of long-term autonomy challenges and sensor evaluation for this purpose is also a significant theme, emphasizing the importance of SLAM technology in real-world applications \cite{b7}. Sensor selection and comparison are crucial, as highlighted in "Sensors, SLAM and Long-term Autonomy: A Review" \cite{b8}, due to their substantial impact on SLAM performance.

The analyzed papers also reveal contrasting viewpoints and approaches, including different SLAM paradigms such as dense and semi-dense methods \cite{b9}, \cite{b10}, as well as deep learning-based approaches like XRDSLAM \cite{b6}. The choice of sensor suites varies across papers, with some focusing on specific sensors like LiDAR \cite{b4} and others discussing various sensor combinations \cite{b8}. Evaluation metrics also differ, encompassing accuracy, energy consumption, processing frame rate, and robustness \cite{b2}, \cite{b3}.

Technical details and methodologies employed in the papers include a range of SLAM algorithms like LOAM, LSD-SLAM, and KinectFusion \cite{b9}, \cite{b10}, \cite{b11}. Semantic segmentation, utilized in SA-LOAM for point cloud data \cite{b4}, has shown promise in enhancing localization accuracy. Modular frameworks and unified dataset management, as seen in XRDSLAM \cite{b6} and GSLAM \cite{b3}, are also significant technical contributions.

Despite the advancements, limitations and challenges persist in SLAM technology. The complexity of evaluating SLAM systems and the lack of benchmarking methodologies are recognized as significant hurdles \cite{b2}, \cite{b3}. Sensor selection and evaluation remain crucial due to their impact on performance \cite{b8}. Scalability and robustness, particularly in large-scale scenes, are also essential considerations, as highlighted by SA-LOAM \cite{b4} and other papers. The transition of SLAM technology from academic research to real-world applications poses a challenge, underscoring the need for long-term autonomy and practical considerations.

In conclusion, the field of SLAM technology is evolving, with a growing emphasis on semantic integration, modularity, flexibility, and long-term autonomy. As highlighted by the analyzed papers, addressing the existing challenges and limitations will be crucial for the prospective future of SLAM. The development of more robust benchmarking frameworks, the exploration of different SLAM paradigms, and the consideration of real-world application requirements will pave the way for significant advancements in SLAM technology \cite{b1}, \cite{b6}, \cite{b8}. Ultimately, these developments will have profound implications for various fields, including robotics, autonomous vehicles, and augmented reality, where accurate and efficient SLAM systems are essential for achieving high-level autonomy and intelligence.

\section{Conclusion}
In conclusion, the analysis of relevant papers on Semantic SLAM reveals significant advancements in the field, underscoring the importance of integrating semantic information into Simultaneous Localization and Mapping (SLAM) systems \cite{b1}. The key findings and contributions from these studies demonstrate that semantic information can enhance localization accuracy, facilitate effective loop closure detection, and enable the construction of globally consistent semantic maps, as evident in approaches like SA-LOAM \cite{b2} and S3LAM \cite{b3}. Furthermore, the use of semantic segmentation to learn robust features for SLAM, as proposed by Semantic SuperPoint \cite{b4}, highlights the potential for improving traditional methods. The development of large-scale datasets, such as DISCOMAN \cite{b5}, has also been crucial for training and benchmarking semantic SLAM methods, emphasizing the need for extensive data resources in this area.

A common theme across these papers is the utilization of semantic information to augment SLAM performance, indicating a shift towards more informed and context-aware mapping and localization techniques. The interdisciplinary nature of Semantic SLAM research is also apparent, with contributions drawing from both computer vision and robotics \cite{b6}. However, differing approaches to incorporating semantic information into SLAM systems are notable, such as the focus on loop closure detection in SA-LOAM versus the emphasis on structured scene understanding in S3LAM \cite{b2}, \cite{b3}. Additionally, novel methods like Semantic SuperPoint offer alternative perspectives on learning robust features for SLAM using semantic segmentation, diverging from conventional techniques \cite{b4}.

The technical methodologies employed in these studies are diverse, encompassing LiDAR-based SLAM, deep learning-based feature extraction, and dataset generation through physically-based rendering, among others \cite{b2}, \cite{b4}, \cite{b5}. The integration of various sensors, including cameras and IMUs, into SLAM systems is also a significant aspect, as discussed in the context of long-term autonomy \cite{b7}. Despite these advancements, challenges persist, such as the need for extensive datasets, the difficulty in detecting loop closures in large-scale environments, and the balance between multiple tasks in multi-task learning approaches \cite{b4}, \cite{b5}. Maintaining long-term autonomy in SLAM systems remains another critical challenge \cite{b7}.

In summary, the surveyed papers collectively demonstrate the potential of semantic information to significantly enhance SLAM performance, while also highlighting key challenges and areas for future research. The integration of semantic insights into SLAM systems is poised to become a cornerstone of advanced mapping and localization technologies, with implications for robotics, autonomous vehicles, and augmented reality applications \cite{b1}, \cite{b6}. As the field continues to evolve, addressing the identified limitations and exploring novel approaches to semantic information integration will be crucial. Future studies should focus on developing more sophisticated methods for incorporating semantic context into SLAM, leveraging advances in computer vision, machine learning, and robotics to create more robust, efficient, and autonomous systems \cite{b2}, \cite{b3}, \cite{b4}.

\end{document}